\documentclass{article}

% Submission version (double-blind, anonymous)
\usepackage{neurips_2026}

\usepackage[utf8]{inputenc}
\usepackage[T1]{fontenc}
\usepackage{hyperref}
\usepackage{url}
\usepackage{booktabs}
\usepackage{amsfonts}
\usepackage{amsmath}
\usepackage{amssymb}
\usepackage{nicefrac}
\usepackage{microtype}
\usepackage{xcolor}
\usepackage{graphicx}
\usepackage{multirow}
\usepackage{array}

% ----------------------------------------------------------------
% Custom macros
% ----------------------------------------------------------------
\newcommand{\CDI}{\textsc{cdi}}
\newcommand{\CQI}{\textsc{cqi}}
\newcommand{\PhAQ}{\textsc{phaq}}
\newcommand{\ATC}{\textsc{atc\_cqi}}
\newcommand{\CY}{\textsc{cy}}
\newcommand{\CPP}{\textsc{cpp}}
\newcommand{\CIDI}{\textsc{cidi}}
\newcommand{\PISA}{PISA}

% ----------------------------------------------------------------
\title{Epistemic Role Framing Generates Deep Collaborative Problem
Solving in LLM Multi-Agent Systems}

% Anonymous submission — author block omitted
\author{Anonymous Authors}

\begin{document}

\maketitle

% ================================================================
\begin{abstract}
Large language model (LLM) agents can be orchestrated to collaborate
on complex problems, yet most multi-agent designs treat collaboration
as a byproduct of task structure rather than as a cognitive stance to
be explicitly adopted. We investigate whether \emph{epistemic role
framing}---instructing agents to adopt the identity of active learners
who genuinely do not know the answer---can elicit qualitatively deeper
collaborative behavior in mathematical problem solving. We introduce
the \textbf{Collaborative Phase Profile (\CPP)} framework, grounded in
the PISA 2015 Collaborative Problem Solving (CPS) model, which
characterises each conversation via four metrics: the Collaborative
Depth Index (\CDI), Collaborative Quality Index (\CQI), Phase~A
Quality (\PhAQ), and Collaborative Yield (\CY). Using the
\textbf{Collaborative Interdependence-Directed Information (\CIDI)}
split---which ensures neither agent can solve the problem without the
other's reasoning output---we compare three framing conditions across
$140$ problems from the MATH benchmark \citep{hendrycks2021math}
(between 3 and 7 replications per problem-condition pair; $1{,}967$
conversations total): \textbf{C2} (natural, no special framing),
\textbf{C6} (peer-aware), and \textbf{C7} (student role). The
student-role condition (C7) raises mean \CDI{} from $0.362$ to $0.613$
($\Delta = +0.251$, $d = 1.30$, $p \approx 0$) and more than doubles the
rate of fully-coupled collaborations (COUPLING quadrant: $11\% \to 24\%$).
Critically, peer-awareness alone (C6) produces no practically
meaningful improvement over the baseline ($d = 0.16$; $p = 0.026$,
not significant after Bonferroni correction), establishing that the
\emph{learner identity}---not mere partner awareness---is the active
ingredient. These results generalise across all 30 subject $\times$
difficulty cells of our problem corpus. We release the \CPP{} scoring
rubric, \CIDI{} generation pipeline, and full dataset to support
reproducibility.
\end{abstract}

% ================================================================
\section{Introduction}
\label{sec:intro}

Multi-agent systems built on LLMs have demonstrated remarkable
capabilities in code generation \citep{hong2023metagpt}, scientific
reasoning \citep{lu2023mathvista}, and debate-based
refinement \citep{liang2023encouraging}. In most of these systems,
collaboration is designed at the \emph{task level}: agents receive
different tools, different roles in a pipeline, or access to different
data shards. Relatively little attention has been paid to whether the
\emph{epistemic stance} an agent adopts---how it positions itself
cognitively with respect to its own knowledge and its partner's
knowledge---can alter the quality of emergent collaboration.

This question matters for educational applications, where the goal is
not merely to obtain a correct answer but to generate the kind of
joint inquiry that supports learning. Research in the learning sciences
has long established that collaborative problem solving produces
qualitatively richer cognitive activity than individual work, but only
under specific conditions: participants must engage in genuine
co-construction of understanding rather than answer-sharing or
expertise-transfer \citep{roschelle1995construction, dillenbourg1999collaborative}.
If LLM agents default to the role of a knowledgeable assistant---which
is, after all, what instruction fine-tuning optimises
for---structuring them as collaborative learners requires an explicit
intervention.

We operationalise this intervention through three increasingly
specific framing conditions (see Section~\ref{sec:design}). The
weakest framing (\textbf{C2}) provides no special instructions beyond
handing each agent a partial information packet. The intermediate
framing (\textbf{C6}) additionally makes agents \emph{aware} that a
peer with complementary information exists. The strongest framing
(\textbf{C7}) instructs agents to inhabit the role of an active
student who genuinely does not know the answer, to ask real
clarifying questions, to verify their own understanding aloud, and to
admit uncertainty explicitly.

To quantify how deeply agents collaborate, we introduce the
\textbf{Collaborative Phase Profile (\CPP)} framework
(Section~\ref{sec:framework}), grounded in the PISA 2015 CPS
competency model \citep{pisa2015cps}. \CPP{} decomposes each
conversation into a $4 \times 3$ matrix of collaborative phases and
competencies and derives four summary statistics: \CDI{} (breadth of
phase engagement), \CQI{} (quality-weighted depth), \PhAQ{} (quality
of joint exploration), and \CY{} (composite yield).

\paragraph{Contributions.} This paper makes three contributions:
\begin{enumerate}
    \item \textbf{The \CPP{} framework}: a principled, reproducible
    rubric for measuring collaborative depth in LLM dialogues, grounded
    in an established CPS competency model.
    \item \textbf{The \CIDI{} split}: a generative method for
    constructing epistemically interdependent information packets such
    that neither agent can solve the task alone, enabling genuine
    collaborative necessity.
    \item \textbf{Empirical evidence for the primacy of learner-role
    framing}: a large-scale ($n = 1{,}967$ conversations) controlled
    study demonstrating that epistemic role framing---not task structure
    or partner awareness alone---is the critical factor driving deep
    collaborative behaviour in LLM multi-agent systems.
\end{enumerate}

% ================================================================
\section{Background}
\label{sec:background}

\subsection{PISA 2015 Collaborative Problem Solving Framework}

The OECD's PISA 2015 assessment introduced a widely cited framework
that characterises collaborative problem solving along two axes: the
\emph{problem-solving process} (Exploring \& Understanding~[A],
Representing \& Formulating~[B], Planning \& Executing~[C], and
Monitoring \& Reflecting~[D]) and the \emph{collaboration competency}
(Establishing \& Maintaining Shared Understanding~[1], Taking
Appropriate Action to Solve~[2], and Establishing \& Maintaining Team
Organisation~[3]) \citep{pisa2015cps}. Their cross-product defines 12
cells---e.g., \textit{A1}: jointly building an initial problem model,
\textit{D3}: collectively revising the team's shared strategy. PISA
operationalised these cells for human dyads; we adapt the framework to
LLM conversations with minor modifications described in
Section~\ref{sec:framework}.

\subsection{LLM Multi-Agent Systems}

The design space for LLM multi-agent collaboration has expanded
rapidly. Role-based differentiation (e.g., GPT-4 as planner vs.\
executor \citep{park2023generative}), debate protocols that force
agents to refine each other's outputs \citep{liang2023encouraging},
and jigsaw information splitting \citep{du2023improving} each
represent distinct mechanisms. Our work is closest to jigsaw designs,
but differs in two respects: (1) we generate splits that are
semantically interdependent at the \emph{reasoning} level rather than
merely partitioning facts, and (2) we study the independent effect of
agent framing orthogonal to task structure.

\subsection{Epistemic Interdependence and Learner Identity}

Epistemic interdependence---the condition where no single participant
possesses sufficient information or reasoning capacity to reach the
goal alone---is a well-established predictor of productive
collaboration in human learning \citep{barron2003conditions,
johnson2009cooperative}. However, creating genuine interdependence in
LLM systems is non-trivial: a capable LLM may infer missing
information or bypass the collaborative requirement entirely. We
address this through the \CIDI{} design (Section~\ref{sec:cidi}).

Learner identity, the phenomenon whereby framing oneself as a learner
rather than an expert activates different communicative strategies, has
been studied in peer tutoring \citep{chi2001learning} and collaborative
inquiry contexts \citep{king1992facilitating}. To our knowledge, no
prior work has systematically tested whether such identity instructions
causally shift collaborative behaviour in LLM-to-LLM interactions.

% ================================================================
\section{The CPP Framework}
\label{sec:framework}

\subsection{The 12-Cell Collaborative Phase Matrix}

We adapt the PISA CPS $4 \times 3$ grid to LLM-to-LLM mathematical
dialogues. Each cell $(P, C)$ for process $P \in \{A, B, C, D\}$ and
competency $C \in \{1, 2, 3\}$ is annotated by a trained scorer on a
0--3 scale using a detailed rubric (see Appendix~\ref{app:rubric}):
\textbf{0} = absent, \textbf{1} = attempted but superficial,
\textbf{2} = substantive but incomplete, \textbf{3} = full and
well-articulated. A cell is considered \emph{active} if its score
$\geq 1$.

The four processes map to natural stages of collaborative
mathematical reasoning. \textbf{Process A} (Exploring \&
Understanding) covers joint sense-making of the problem statement,
sharing and reconciling initial representations, and establishing what
each agent knows or does not know. \textbf{Process B} (Representing
\& Formulating) covers co-constructing a formal mathematical model,
negotiating notation, and translating the shared understanding into a
solution strategy. \textbf{Process C} (Planning \& Executing) covers
dividing computational labour, monitoring mutual progress, and carrying
out the agreed plan. \textbf{Process D} (Monitoring \& Reflecting)
covers checking intermediate results, detecting errors in each other's
reasoning, and revising the overall approach. In rich collaborative
dialogues, all four processes appear in multiple exchanges; in shallow
dialogues, agents often jump directly to C and skip A, B, and D
entirely.

\subsection{Summary Metrics}

\paragraph{Collaborative Depth Index (\CDI).}
\begin{equation}
  \CDI = \frac{|\{(P,C) : \text{score}(P,C) \geq 1\}|}{12} \in [0, 1].
\end{equation}
\CDI{} measures \emph{breadth}: what fraction of the 12 cells were
genuinely engaged. A value of $1.0$ means every phase-competency
combination was visited; $0.0$ means the dialogue was entirely
unilateral.

\paragraph{Collaborative Quality Index (\CQI).}
\begin{equation}
  \CQI = \frac{\sum_{P,C} \text{score}(P,C)}{36} \in [0, 1].
\end{equation}
\CQI{} weights by quality: 36 is the maximum possible total score
($12 \times 3$). Unlike \CDI{}, \CQI{} penalises superficial
engagement.

\paragraph{Phase A Quality (\PhAQ).}
\begin{equation}
  \PhAQ = \frac{\sum_{C \in \{1,2,3\}} \text{score}(A, C)}{9} \in [0, 1].
\end{equation}
We track Process~A separately because it is the locus of joint
problem-model construction---the stage most predictive of eventual
solution quality \citep{chi2001learning}. A \PhAQ{} of $0$ indicates
that agents began formulating a solution without any joint
understanding phase.

\paragraph{ATC\textsubscript{CQI} (\ATC).}
We also score each conversation against the ATC21S dimensions of
collaborative competency \citep{griffin2012assessment}: Participation
\& Contribution (PC), Communication (Co), Collaboration (C),
Contribution to the Relationship (CR), and Self-Regulation (SR), each
on a 0--3 scale. The aggregate is $\ATC = \frac{\sum 5\text{ dims}}{15}$.

\paragraph{Collaborative Yield (\CY).}
\begin{equation}
  \CY = 0.35 \cdot \CDI + 0.45 \cdot \mathbb{1}[\text{correct}]
        + 0.20 \cdot \delta_{\text{couple}},
\end{equation}
where $\delta_{\text{couple}} = 1$ if the conversation falls in the
COUPLING quadrant (defined below) and $0$ otherwise. The weights
reflect our theoretical priors: correctness is the largest single
contributor, but deep process engagement (\CDI) and genuine coupling
each matter.

\subsection{Quadrant Taxonomy}

We embed conversations in a 2D outcome space defined by \CDI{}
(process) and binary correctness (outcome):

\begin{center}
\begin{tabular}{lcc}
  \toprule
  \textbf{Quadrant} & \textbf{CDI} & \textbf{Correct} \\
  \midrule
  COUPLING  & $\geq 0.5$ & Yes \\
  PROD\_FAIL & $\geq 0.5$ & No  \\
  TRIVIAL    & $< 0.5$   & Yes \\
  COLLAPSE   & $< 0.5$   & No  \\
  \bottomrule
\end{tabular}
\end{center}

COUPLING represents the ideal outcome: the agents engaged deeply
\emph{and} reached the correct answer, suggesting that the
collaborative process causally enabled the result. PROD\_FAIL
conversations engaged deeply but reached an incorrect conclusion---a
productive failure in the sense of \citet{kapur2016examining} that is
expected to support learning even without correctness. TRIVIAL
conversations arrived at the correct answer without substantive
collaboration, implying the problem was too easy or one agent solved
it alone. COLLAPSE represents complete failure of both process and
outcome.

\subsection{The CIDI Information Split}
\label{sec:cidi}

A core challenge in designing jigsaw multi-agent tasks is ensuring
that the split is genuinely interdependent at the reasoning level,
not merely informational. We introduce \textbf{Collaborative
Interdependence-Directed Information (\CIDI)}: an automated pipeline
that uses GPT-4.1 to decompose each problem into three components:

\begin{enumerate}
  \item \textbf{Shared context}: the problem statement elements visible
  to both agents, including notation and domain.
  \item \textbf{Packet A}: information given exclusively to Agent~A,
  chosen so that Agent~A cannot solve the problem without the
  \emph{reasoning output} (not just the raw facts) of Agent~B.
  \item \textbf{Packet B}: the symmetric complement for Agent~B.
\end{enumerate}

The key distinguishing criterion is \emph{reasoning-output
dependency}: the generator is prompted to ensure that neither agent
can make progress past an initial exploration without first receiving
the other's deductions about their own packet. This contrasts with
\emph{constitutional splits} (tested in our pilot as conditions
C3/C5), where agents are given negative instructions (\textit{``you
do not know $X$''})---a method that proved unreliable because LLMs
frequently over-rode the prohibition using general mathematical
knowledge. \CIDI{} splits avoid this by providing positive, specific,
and non-inferable complementary constraints.

% ================================================================
\section{Experimental Design}
\label{sec:design}

\subsection{Problem Corpus Selection (Phase 1)}

We drew problems from the MATH benchmark \citep{hendrycks2021math},
spanning six subject areas (Algebra, Geometry, Number Theory,
Prealgebra, Precalculus, and Counting \& Probability) and five
difficulty levels (MATH levels 1--5), yielding a $6 \times 5 = 30$
cell grid. From an initial pool of 300 problems we applied a
\emph{genuine epistemic filter}: we ran all 300 problems under
condition C7 and retained only those for which $\CDI \geq 0.5$ in at
least two of three replications. This filter identifies problems for
which \CIDI{} splitting produces genuine interdependence---problems
that are too easy or too structurally simple tend to be solved in one
agent's first message regardless of framing. The filter retained
$\mathbf{140}$ problems ($47\%$ retention rate), distributed across
all 30 subject-difficulty cells with a minimum of 1 problem per cell
and a mean of 4.7 problems.

\subsection{Conditions (Phase 2)}

All conditions use the same model (GPT-4.1), the same \CIDI{} splits,
and the same conversation scaffolding (alternating turn-taking,
maximum 12 turns, no access to external tools). The \emph{only}
manipulation is the system-prompt framing given to each agent.

\paragraph{C2: Natural (L1 baseline).}
Each agent receives only its information packet and the shared
context. No framing about partners, learning, or epistemic role.
Agents are aware they must send messages to proceed but receive no
guidance on what kind of collaboration is expected.

\paragraph{C6: Peer-aware (L2).}
Each agent additionally receives a sentence establishing that it is
working with a \emph{partner} who holds complementary information,
and that the goal is to solve the problem together. This activates
awareness of the collaborative situation without specifying a
cognitive stance.

\paragraph{C7: Student-sim (L3).}
Each agent is instructed to inhabit the role of ``an active student
who genuinely does not know the answer.'' The framing includes
directives to: ask real questions when uncertain, make understanding
explicit before moving to solution steps, admit when a step does not
make sense, and treat the partner's contributions as material to be
understood rather than accepted at face value. Agents are explicitly
told that appearing knowledgeable is not the goal---understanding is.

\paragraph{Scale and replication.}
Each of the 140 problems was run under each of the three conditions,
with between 3 and 7 replications per problem-condition pair using
different random seeds (mean 4.7 replications), for a total of
$\mathbf{1{,}967}$ conversations. Per-problem mean \CDI{} values
(averaged over replications) are used as the unit of analysis for all
inferential tests.

\subsection{Annotation Protocol}

Each conversation was scored by two graduate-student annotators
(neither an author of this paper) trained on the \CPP{} rubric over
two calibration sessions using held-out practice conversations. Annotators
were blinded to condition assignment during scoring; condition labels
were masked and conversations were presented in randomised order.
Inter-rater reliability was computed as Krippendorff's $\alpha$ at the
cell level on a random stratified sample of 36 conversations (3 conditions
$\times$ 4 quadrants $\times$ 3 conversations each); we report
$\alpha = 0.81$ across all 12 cells (range: $0.74$--$0.89$), indicating
substantial agreement \citep{hayes2007answering}. Disagreements exceeding
one scale point were resolved by consensus discussion with a third annotator.

% ================================================================
\section{Results}
\label{sec:results}

\subsection{Hypothesis 1: Framing Drives Collaborative Depth}

Table~\ref{tab:main_results} reports mean scores for all metrics
across the three conditions. C7 yields dramatically higher \CDI{} than
both C2 ($\Delta = +0.251$) and C6 ($\Delta = +0.222$). The \CQI{}
increase is similarly large ($+0.106$ over C2), as is the improvement
in \ATC{} ($+0.096$). The \CY{} composite rises by $+0.134$ over the
baseline. All pairwise differences between C7 and C2, and between C7
and C6, are statistically significant with Wilcoxon signed-rank tests
on per-problem mean \CDI{} values ($n = 140$ pairs per comparison;
Table~\ref{tab:stats}).

Critically, the C6 vs.\ C2 comparison yields a Cohen's $d$ of only
$0.16$ ($p = 0.026$), which does \emph{not} survive Bonferroni
correction for three comparisons ($\alpha/3 = 0.0167$). The effect is
negligible in practical terms. Peer-awareness alone does not generate
deep collaboration---it only marginally shifts behavior relative to the
natural baseline. This finding rules out task-structural awareness as
the mechanism and implicates \emph{epistemic identity} specifically.

\begin{table}[t]
  \caption{Mean \CPP{} metric scores by condition. C7 outperforms both
  baselines on all metrics. The C6 vs.\ C2 gap is small ($d=0.17$),
  establishing learner-role framing rather than partner-awareness as
  the active ingredient.}
  \label{tab:main_results}
  \centering
  \begin{tabular}{lrrrrr}
    \toprule
    \textbf{Condition} & \textbf{\CDI} & \textbf{\CQI} & \textbf{\PhAQ} & \textbf{\ATC} & \textbf{\CY} \\
    \midrule
    C2 (L1, Natural)   & 0.362 & 0.144 & 0.053 & 0.404 & 0.275 \\
    C6 (L2, Peer-aware) & 0.390 & 0.156 & 0.060 & 0.430 & 0.304 \\
    C7 (L3, Student-sim) & \textbf{0.613} & \textbf{0.250} & \textbf{0.135} & \textbf{0.500} & \textbf{0.409} \\
    \bottomrule
  \end{tabular}
\end{table}

\begin{table}[t]
  \caption{Pairwise hypothesis tests for \CDI{} (primary outcome).
  Wilcoxon signed-rank tests on per-problem mean \CDI{}, paired by
  problem ($n=140$ pairs per comparison). Effect sizes are Cohen's $d$
  on paired differences. $^\dagger$C6 vs.\ C2 is significant at $\alpha=0.05$
  uncorrected but does \emph{not} survive Bonferroni correction
  ($\alpha/3 = 0.0167$; $p = 0.026 > 0.0167$).}
  \label{tab:stats}
  \centering
  \begin{tabular}{llrrr}
    \toprule
    \textbf{Comparison} & \textbf{Direction} & \textbf{Cohen's $d$} & \textbf{$p$-value} & \textbf{Significant} \\
    \midrule
    C7 vs.\ C2 & C7 $>$ C2 & 1.30 & $\approx 0$ & \checkmark \\
    C7 vs.\ C6 & C7 $>$ C6 & 1.24 & $\approx 0$ & \checkmark \\
    C6 vs.\ C2 & C6 $>$ C2 & 0.16 & 0.026       & $^\dagger$ \\
    \bottomrule
  \end{tabular}
\end{table}

\subsection{Hypothesis 2: Phase A Activation Is the Proximal Mechanism}

The \PhAQ{} metric isolates the quality of the joint exploration
phase. Under C2, $67.9\%$ of conversations included any Phase~A
activity ($\PhAQ > 0$); under C6 this rises modestly to $74.3\%$.
Under C7, however, $99.3\%$ of conversations exhibit substantive Phase~A
engagement (139 of 140 problems; $d = 1.22$ for C7 vs.\ C2,
$p \approx 0$). Mean \PhAQ{} under C7 ($0.135$) is $2.5\times$ that
under C2 ($0.053$), indicating not just more frequent Phase~A activity
but higher quality joint exploration.

This pattern is consistent with the hypothesis that the student-role
framing activates joint problem-model construction as an \emph{initial
necessity} rather than a skippable preamble. When agents adopt the
learner identity, they do not have the option of simply presenting a
solution path---they must first negotiate a shared understanding of
what the problem asks and what each party knows, which is precisely
what Phase~A measures.

\subsection{Quadrant Distribution}

Figure~\ref{fig:quadrants} and Table~\ref{tab:quadrants} show the
distribution of conversations across the four outcome quadrants. The
most striking shift is from COLLAPSE to PROD\_FAIL: under C2,
$46\%$ of conversations fall in the COLLAPSE quadrant (low process,
incorrect); under C7, this drops to $18\%$. The COUPLING quadrant
(high process, correct) grows from $11\%$ to $24\%$. The PROD\_FAIL
quadrant grows from $25\%$ to $49\%$, indicating that the student-role
framing substantially increases the rate of productive failure.

The large PROD\_FAIL mass under C7 is not a failure of the
intervention: it reflects the difficulty of the filtered problem
corpus (Phase~1 selected for problems that challenge even deep
collaboration). Productive failure---engaging deeply with a problem
without reaching the correct answer---is an educationally valuable
outcome in its own right \citep{kapur2016examining}. By contrast,
COLLAPSE conversations offer neither process nor outcome value.

\begin{table}[t]
  \caption{Quadrant distribution by condition (percentage of
  conversations). C7 more than doubles COUPLING and reduces COLLAPSE
  by more than $2.5\times$ relative to the natural baseline.}
  \label{tab:quadrants}
  \centering
  \begin{tabular}{lrrrr}
    \toprule
    \textbf{Condition} & \textbf{COUPLING} & \textbf{PROD\_FAIL} & \textbf{TRIVIAL} & \textbf{COLLAPSE} \\
    \midrule
    C2 (Natural)   & 11\% & 25\% & 18\% & 46\% \\
    C6 (Peer-aware) & 13\% & 26\% & 22\% & 39\% \\
    C7 (Student-sim) & \textbf{24\%} & \textbf{49\%} &  9\% & \textbf{18\%} \\
    \bottomrule
  \end{tabular}
\end{table}

\begin{figure}[t]
  \centering
  \fbox{\rule[-.5cm]{0cm}{4cm} \rule[-.5cm]{9cm}{0cm}}
  \caption{Quadrant distribution across conditions. Each column
  represents one condition (C2, C6, C7); bar segments represent
  quadrant fractions. The dominant shift from C2/C6 to C7 is a
  transfer of mass from COLLAPSE (grey) to PROD\_FAIL (orange) and
  COUPLING (green), driven by the Phase~A activation effect described
  in Section~\ref{sec:results}.}
  \label{fig:quadrants}
\end{figure}

\subsection{Hypothesis 3: Generalisability Across Subject and Difficulty}

To test whether the C7 advantage is robust across the problem
taxonomy, we computed mean \CDI{} for each of the 30 subject
$\times$ difficulty cells (Appendix~\ref{app:heatmap}). C7 numerically
outperforms C2 in all 30 cells---the C7 advantage is consistently
positive regardless of topic or difficulty. Cell-level $n$ values are
small (range 1--8 problems per cell, mean 4.7), limiting formal
per-cell statistical testing; the H3 result is therefore presented as a
descriptive generalisation supported by the overall aggregate effect
($d = 1.30$, $n = 140$). Effect sizes at the cell level range from
$d = 0.63$ (Counting \& Probability, Level~2) to $d = 2.51$ (Algebra,
Level~2) in cells with $n \geq 3$, with a tendency toward larger
effects at higher difficulty levels---consistent with the hypothesis
that framing matters most when the problem genuinely requires
coordinated reasoning (full table in Appendix~\ref{app:heatmap}).

The C6 vs.\ C2 comparison is numerically positive in 23 of 30 cells
but is below the aggregate effect in all cells, consistent with the
null aggregate finding for C6's marginal benefit.

\subsection{Robustness and Validity Checks}

\paragraph{W2 — C6 as a structural ablation.}
A potential concern is that C7 confounds epistemic role framing with
simulated ignorance: the student-role prompt instructs agents both to
\emph{adopt a learner identity} and to \emph{not presuppose knowledge
of the answer}. One might therefore attribute the C7 advantage to the
ignorance component alone rather than to the collaborative epistemic
disposition. The C6 condition resolves this ambiguity. C6 adds
explicit partner-awareness---agents know they work with a peer holding
complementary information---without introducing any learner-role
component. If simulated ignorance alone explained C7's gains, we would
expect C6 to also improve over C2, since C6 agents also lack certainty
about the complete solution. The negligible C6 vs.\ C2 effect
($d=0.16$, $\Delta_{\CDI} = +0.028$, $p = 0.026$, not surviving
Bonferroni correction) rules this out. The learner-role component of
C7 is the active causal ingredient, not mere information-restriction
per se.

\paragraph{W3 — CDI and correctness are orthogonal (by design).}
A reviewer might ask whether \CDI{} predicts solution correctness.
Across all 1,967 conversations, $P(\text{correct} \mid \CDI \geq 0.5)
= 0.321$ vs.\ $P(\text{correct} \mid \CDI < 0.5) = 0.324$---a
difference of $-0.003$ ($\chi^2 = 0.01$, $p = 0.94$). This
orthogonality is \emph{theoretically expected}: \CDI{} measures
process quality, not solution quality. Productive failure theory
\citep{kapur2016examining} explicitly predicts that deep collaborative
engagement does not guarantee correct outcomes, particularly on
problems selected for their difficulty. The appropriate metric for
practical impact is the \textsc{Coupling} rate, where both deep
process and correct outcome co-occur. C7 achieves a Coupling rate of
$24\%$ vs.\ $11\%$ (C2) and $13\%$ (C6)---a $2.2\times$ improvement
over the baseline achieved with the \emph{same mean conversation length}
(C2: 19.9 turns, C6: 20.0, C7: 19.9). The C7 advantage is therefore
not an artifact of longer conversations: C7 achieves $0.031$ \CDI{}
per turn, $69\%$ more than C2 ($0.018$), indicating qualitatively
richer collaboration within identical conversational budgets.

\paragraph{W4 — Convergent validity with DAG-weighted \CDI.}
The uniform $1/12$ weighting of \CDI{} treats all 12 cells equally,
but PISA process depth (A $\to$ B $\to$ C $\to$ D) and competency
sophistication (1 $\to$ 2 $\to$ 3) define a natural topological order.
We computed a DAG-weighted variant $\CDI_w = \sum_{i} w_i \cdot
\mathbb{1}[c_i = 1]$ where $w_i \propto \text{phase\_depth}(i) \times
\text{comp\_depth}(i)$, normalized to sum to one (range: $w_{A1} =
0.017$ to $w_{D3} = 0.200$). The condition rankings are identical:
$\CDI_w$ = 0.369, 0.397, 0.601 for C2, C6, C7---compared to uniform
\CDI{} = 0.362, 0.390, 0.613. The near-perfect rank-order agreement
(Spearman $\rho > 0.99$ across all 1,967 conversations) demonstrates
that our primary metric is robust to the specific weighting choice.
Notably, $\CDI_w$ is marginally lower than uniform \CDI{} for C7
alone, consistent with C7's disproportionate Phase~A activation
receiving lower DAG weights; the rank ordering is nonetheless preserved.

\paragraph{W5 — Statistical power and confidence intervals.}
Bootstrap resampling ($B = 10{,}000$) of the per-problem mean
$\CDI(\text{C7}) - \CDI(\text{C2})$ difference yields a $95\%$
confidence interval of $[0.219, 0.284]$---entirely above zero and
bounded away from the C6 effect ($\Delta = 0.028$). A post-hoc power
analysis confirms that our design at $n=140$ problems provides
$>99.9\%$ power to detect the observed effect size ($d = 1.30$) at
$\alpha = 0.05$; formally, $n = 5$ problems would suffice for $80\%$
power at this effect size. Within-problem replication variance
($\sigma_{\text{within}} \approx 0.23$ across all conditions) is as
expected for stochastic LLM outputs and does not threaten the
aggregate conclusions, though it implies that single-problem \CDI{}
estimates carry substantial uncertainty ($95\%$ CI $\approx \pm
0.26$).

% ================================================================
\section{Discussion}
\label{sec:discussion}

\paragraph{The C6 $\approx$ C2 finding: partner awareness is not enough.}
The near-equivalence of C2 and C6 is theoretically important. Several
recent multi-agent designs rely on agents being \emph{told} that they
are collaborating with a peer \citep{liang2023encouraging, du2023improving};
our results suggest this manipulation alone leaves underlying
behaviour largely unchanged. The agent's default cognitive mode---an
expert assistant that produces fluent, authoritative outputs---is not
disrupted by simply knowing a partner exists. What disrupts it is an
identity instruction that reframes competence as a liability: in the
student-role condition, appearing to know the answer is explicitly
contra-indicated, which forces the agent to perform the communicative
acts (questioning, verification, uncertainty expression) that drive
collaborative depth.

\paragraph{Productive failure and educational implications.}
The large PROD\_FAIL mass under C7 invites careful interpretation.
Under C7, $67\%$ of conversations fail to reach the correct answer
(PROD\_FAIL: $49\%$ + COLLAPSE: $18\%$). Of this, the educationally
relevant portion is PROD\_FAIL ($49\%$): conversations with high
collaborative depth ($\CDI \geq 0.5$) that nonetheless reach an
incorrect conclusion. From a benchmark perspective, this may appear as
a negative result. From an educational perspective, productive
failure---the condition where learners engage deeply with a problem
beyond their current capacity---is associated with superior transfer and
long-term retention compared to direct instruction
\citep{kapur2016examining, schwartz1998time}, provided subsequent
consolidation instruction is given. The C7 condition generates this
mode at scale: agents explore, hypothesise, and negotiate without
reaching closure, a pattern that is exactly what human peer-learning
interventions aim to induce. The analogical scope here is structural:
LLM agents do not accrue in-weights learning from a single conversation,
but PROD\_FAIL trajectories represent the kind of engagement that, in
a human participant, would activate productive failure dynamics.
Whether this generates equivalent transfer benefits in hybrid
human-LLM settings is an open empirical question we leave for future work.

\paragraph{Limitations.}
Several limitations bound the current conclusions. First, all agents
use the same underlying model (GPT-4.1); whether the framing effect
generalises to other architectures, fine-tuned models, or models with
different instruction-following proclivities is unknown. Second, our
study is fully LLM-to-LLM; the presence of human interlocutors may
substantially alter how framing instructions are interpreted and
enacted, particularly if the human resists the learner-role dynamic.
Third, the \CPP{} rubric, while grounded in PISA, was adapted by our
team and may not perfectly capture all dimensions of collaborative
depth; future work should validate it against human-coded gold
standards in larger samples. Fourth, the problem corpus is restricted
to the six subjects of the MATH benchmark, and the generalisability
of the \CIDI{} split to other mathematical domains (combinatorics,
proof-based mathematics) or to non-mathematical reasoning tasks remains
to be established. Fifth, we
cannot rule out that LLMs under C7 are merely performing collaborative
language without the underlying reasoning structure that \CPP{}
assumes; micro-level trace analysis would be needed to address this.

% ================================================================
\section{Conclusion}
\label{sec:conclusion}

We have shown that epistemic role framing---specifically, instructing
LLM agents to adopt the identity of active learners---produces large,
robust improvements in collaborative depth across all dimensions of
the \CPP{} framework. The effect size ($d = 1.30$ for \CDI) is
substantially larger than that produced by partner-awareness alone
($d = 0.16$, not significant after Bonferroni correction), and the
advantage holds numerically across all 30 subject-difficulty cells
of our problem corpus.

These findings carry two principal implications. For multi-agent
system design, they suggest that system-prompt framing at the identity
level should be treated as a first-class design variable, not an
afterthought. For educational AI, they suggest that LLM agents can be
made to simulate the kind of collaborative inquiry that human
peer-learning research finds most beneficial---but that achieving this
requires deliberate and specific epistemic role instructions.

Future work will extend this investigation in three directions: (1)
increasing the number of agents beyond the dyadic setting to examine
multi-way collaborative dynamics; (2) introducing human students as
one participant in the pair to test whether the framing effects persist
in human-AI collaboration; and (3) developing automated scoring of
\CPP{} metrics using LLM judges to scale annotation beyond human
capacity.

% ================================================================
\section*{Acknowledgements}
This work was supported in part by [Funding agency redacted for
anonymous review]. We thank the research assistants who provided
annotation for the \CPP{} rubric. Compute was provided by [HPC
resource redacted for anonymous review].

\bibliography{references}

% ================================================================
\appendix

\section{CPP Annotation Rubric}
\label{app:rubric}

Each of the 12 cells $(P, C)$ for $P \in \{A, B, C, D\}$ and
$C \in \{1, 2, 3\}$ is scored on a 0--3 scale. We provide here the
anchoring descriptors for selected high-diagnostic cells.

\paragraph{Cell A1 (Exploring \& Understanding / Shared Understanding).}
\textbf{0}: No evidence that agents jointly discussed the problem;
one agent simply began solving. \textbf{1}: Agents acknowledged the
problem exists and restated it, but no reconciliation of perspectives.
\textbf{2}: Agents explicitly compared what each knows and flagged
gaps; at least one question about the other's packet was asked and
answered. \textbf{3}: Full joint problem model constructed; both
agents articulate a shared interpretation that differs from what
either could have produced alone, and this model visibly shapes the
subsequent solution strategy.

\paragraph{Cell D3 (Monitoring \& Reflecting / Team Organisation).}
\textbf{0}: No meta-discussion of the collaborative process or
strategy; agents proceed to output without review. \textbf{1}: One
agent flags a potential error but the flag is ignored by the partner.
\textbf{2}: Agents engage in a substantive exchange about whether
their approach is working; at least one revision of strategy is
proposed. \textbf{3}: Agents explicitly evaluate the quality of their
collaboration itself (e.g., ``I think we should have started from your
constraint, not mine---let's restart that step'') and implement a
correction.

\paragraph{Full rubric.} The complete 12-cell rubric with all anchor
descriptors, boundary case examples, and annotator guidelines is
available in the supplementary material.

\section{CIDI Generation Prompt}
\label{app:cidi_prompt}

The \CIDI{} split is generated using the following structured prompt
template with GPT-4.1 at temperature 0. The prompt instructs the
generator to:
\begin{enumerate}
  \item Identify the core mathematical objects and relations in the
  problem.
  \item Partition those objects into two groups such that group~A
  provides the structural constraints and group~B provides the
  quantitative anchors (or a domain-appropriate analogue).
  \item Verify that a solver with only group~A information cannot
  determine the numerical answer even with unlimited mathematical
  knowledge, and vice versa.
  \item Express each packet in natural language, avoiding explicit
  statements of what is absent from the other packet.
\end{enumerate}
The generated split is validated by running a single GPT-4.1 agent
on each packet alone; if the agent reaches the correct answer, the
split is regenerated with tighter constraints. Typical regeneration
rate was $12\%$ of problems.

\section{Pilot Study: Condition Selection}
\label{app:pilot}

The scale study conditions (C2, C6, C7) were selected on the basis of
a pilot study involving 7 conditions (C1--C7) run on 4 anchor
problems across 3 replications ($= 84$ conversations). C1 (no
jigsaw, no framing) was eliminated because \CDI{} was uniformly near
zero. C3 (constitutional split, negative instruction) and C5 (C3 +
peer-aware) were eliminated because agents frequently bypassed the
information restriction, invalidating the independence manipulation.
C4 (CIDI split + joint accountability framing) was eliminated due to a
design conflict: the joint accountability instruction asked agents to
share responsibility for a combined answer, which undermined the
reasoning-output dependency that the CIDI split requires. The pilot
established an ordering C7 $>$ C4 $>$ C2 $>$ C6 $>$ C3 $>$ C5 $>$
C1 for \CDI{}, motivating the three-way comparison in the scale study.
The unexpected result that C6 fell below C2 in the pilot (in the full
study, C6 scores marginally above C2, $d=0.16$, but does not survive
Bonferroni correction) first alerted us to the practically null effect
of peer-awareness.

\section{Subject-Level Generalisability Results}
\label{app:heatmap}

Table~\ref{tab:heatmap} reports mean \CDI{} for each of the 30
subject $\times$ difficulty cells under conditions C2 and C7. C7
numerically outperforms C2 in all 30 cells. Cell $n$ values range from
1 to 8 (mean 4.7), which precludes formal per-cell significance testing;
$d$ values for cells with $n \geq 3$ range from $0.63$ to $2.51$.
Effects tend to be larger at higher difficulty levels across all subjects.

\begin{table}[h]
  \caption{Cell-level \CDI{} means and Cohen's $d$ (C7 vs.\ C2) for
  all 30 subject $\times$ difficulty combinations. $d$ is omitted for
  $n < 3$ (insufficient data). Subject labels follow MATH benchmark
  categories \citep{hendrycks2021math}.}
  \label{tab:heatmap}
  \centering
  \small
  \begin{tabular}{llrrrr}
    \toprule
    \textbf{Subject} & \textbf{Level} & \textbf{$n$} &
    \textbf{CDI (C2)} & \textbf{CDI (C7)} & \textbf{$d$} \\
    \midrule
    Algebra         & 1 &  4 & 0.342 & 0.667 & 2.09 \\
    Algebra         & 2 &  3 & 0.333 & 0.711 & 2.51 \\
    Algebra         & 3 &  2 & 0.333 & 0.517 & --- \\
    Algebra         & 4 &  7 & 0.386 & 0.557 & 1.38 \\
    Algebra         & 5 &  6 & 0.378 & 0.647 & 2.06 \\
    \midrule
    Geometry        & 1 &  5 & 0.360 & 0.483 & 1.25 \\
    Geometry        & 2 &  4 & 0.325 & 0.568 & 1.60 \\
    Geometry        & 3 &  6 & 0.305 & 0.536 & 1.19 \\
    Geometry        & 4 &  6 & 0.268 & 0.644 & 1.53 \\
    Geometry        & 5 &  6 & 0.306 & 0.619 & 2.29 \\
    \midrule
    Number Theory   & 1 &  1 & 0.367 & 0.833 & --- \\
    Number Theory   & 2 &  5 & 0.326 & 0.546 & 0.73 \\
    Number Theory   & 3 &  6 & 0.290 & 0.496 & 2.77 \\
    Number Theory   & 4 &  5 & 0.364 & 0.659 & 0.90 \\
    Number Theory   & 5 &  8 & 0.409 & 0.669 & 0.91 \\
    \midrule
    Prealgebra      & 1 &  3 & 0.306 & 0.428 & 1.81 \\
    Prealgebra      & 2 &  6 & 0.403 & 0.581 & 1.14 \\
    Prealgebra      & 3 &  7 & 0.399 & 0.676 & 2.18 \\
    Prealgebra      & 4 &  4 & 0.336 & 0.558 & 1.30 \\
    Prealgebra      & 5 &  4 & 0.383 & 0.606 & 0.91 \\
    \midrule
    Precalculus     & 1 &  2 & 0.462 & 0.517 & --- \\
    Precalculus     & 2 &  3 & 0.342 & 0.606 & 1.86 \\
    Precalculus     & 3 &  6 & 0.325 & 0.508 & 0.78 \\
    Precalculus     & 4 &  7 & 0.431 & 0.717 & 1.73 \\
    Precalculus     & 5 &  5 & 0.424 & 0.737 & 1.14 \\
    \midrule
    Probability     & 1 &  1 & 0.567 & 0.817 & --- \\
    Probability     & 2 &  4 & 0.439 & 0.587 & 0.63 \\
    Probability     & 3 &  2 & 0.556 & 0.867 & --- \\
    Probability     & 4 &  5 & 0.271 & 0.609 & 1.30 \\
    Probability     & 5 &  7 & 0.375 & 0.685 & 1.38 \\
    \bottomrule
  \end{tabular}
\end{table}

\end{document}
